\begin{definition}[The admissible cone $\mathcal{K}$]
With $\mathcal{V} = \{v \in L^2(\gamma_n) : \mathcal{E}(v) < \infty\}$ the form
domain of \noderef{Energy}, set
\[
  \mathcal{K} = \{v \in \mathcal{V} : v \geq 0 \ \ \gamma_n\text{-a.e.}\} .
\]
This is the admissible set of the obstacle problem: the nonnegative elements of
the fractional form domain.  It is nonempty, since $0 \in \mathcal{K}$, and it is
a convex cone: if $v, w \in \mathcal{K}$ and $\lambda, \mu \geq 0$ then
$\lambda v + \mu w \geq 0$ a.e., and $\mathcal{E}(\lambda v + \mu w) < \infty$
because $\mathcal{V}$ is a linear subspace of $L^2(\gamma_n)$.  Constants
$c \geq 0$ belong to $\mathcal{K}$, with $\mathcal{E}(c) = 0$.
\end{definition}
